% Keep styling.sty beside this file. Compile with:
% latexmk -xelatex -synctex=1 template.tex (or, 
% for my purposes, compile to preview with VimtexCompile in neovim)
% if youre using VSC with some built in plugin you'll probably need to add 
% {
%   "latex-workshop.latex.recipe.default": "latexmk (xelatex)",
%   "latex-workshop.latex.autoBuild.run": "onSave",
%   "latex-workshop.view.pdf.viewer": "tab",
%   "latex-workshop.view.pdf.invertMode.enabled": "never"
% }
% also it is necessary to download the right fonts. Just look 
% at error messages to find which ones or something

\documentclass[11pt]{article}
\usepackage{styling}
% \usepackage{graphicx} % Uncomment if adding external images.



\usepackage[table]{xcolor}

\providecommand{\GridYes}[1]{\cellcolor{green!30}\textcolor{black}{#1}}
\providecommand{\GridNo}[1]{\cellcolor{red!27}\textcolor{black}{#1}}
\providecommand{\GridWork}[1]{\cellcolor{yellow!25}\textcolor{black}{#1}}
\providecommand{\GridNA}{\cellcolor{black!22}\textcolor{black}{\textsf{N/A}}}
\providecommand{\GridHint}[1]{\textit{#1}}
\providecommand{\GridHead}[1]{\rotatebox{90}{\strut #1}}


\title{Hand in 2}
\author{Thobias K. H\o ivik}
\date{For October 23, 2026} % Use \date{} to hide the date.

\begin{document}
\maketitle
% \tableofcontents
% \newpage

Sorry for the miniuscule text, the tables overflowed 
when I tried to make it bigger. 
Also, the links to propositions and such 
do not provide much contrast to the cell colors, but 
they should be clickable in most PDF readers. 

\begingroup
\setlength{\tabcolsep}{5pt}
\renewcommand{\arraystretch}{1.45}
\begin{center}
\textbf{Exercise 1(d--j): Properties of particular spaces}

\medskip
\resizebox{\linewidth}{!}{%
\begin{tabular}{|l|*{7}{c|}}
\hline
 & $I_o^2$ & $\mathbb R_\ell$ & $\mathbb R_\ell^2$
 & $\mathbb R^\omega$ product & $\mathbb R^\omega$ uniform
 & $\mathbb R^\omega$ box & $\mathbb R^{[0,1]}$ product \\
\hline
Connected
 & \GridYes{Example 6, \S24} & \GridNo{$[0,\infty)$}
 & \GridNo{$[0,\infty)\times\mathbb R_\ell$}
 & \GridYes{Exercise 10, \S23} & \GridNo{See \cref{prop:omega_uniform_connectedness}}
 & \GridNo{Example 6, \S23} & \GridYes{Exercise 10, \S23} \\
\hline
Path connected
 & \GridNo{Example 6, \S24} & \GridNo{$\uparrow$}
 & \GridNo{$\uparrow$} & \GridYes{Interpolation}
 & \GridWork{\GridNo{Implied by above}} & \GridNo{$\uparrow$}
 & \GridWork{\GridHint{}} \\
\hline
Locally connected
 & \GridYes{Exercise 3, \S25} & \GridNo{$[y,\infty)$}
 & \GridNo{$\prod_i[y_i,\infty)$} & \GridYes{$\downarrow$}
 & \GridYes{See \cref{prop:omega_uniform_local_connectedness}} & \GridNo{\GridHint{Implied by below}}
 & \GridWork{\GridHint{}} \\
\hline
Locally path connected
 & \GridNo{Exercise 3, \S25} & \GridNo{$\uparrow$}
 & \GridNo{$\uparrow$} & \GridYes{Interpolation}
 & \GridYes{\GridHint{See \cref{prop:omega_uniform_local_connectedness}}} & \GridNo{\GridHint{See \cref{prop:box_locally_path_connected}}}
 & \GridWork{\GridHint{}} \\
\hline
Compact
 & \GridWork{\GridYes{Thm 27.1}} & \GridNo{$\{[n,n+1)\}_{n\in\mathbb Z}$}
 & \GridNo{$\leftarrow$} & \GridNo{$\pi_i$ onto $\mathbb R$}
 & \GridNo{$\pi_i$ onto $\mathbb R$} & \GridNo{$\pi_i$ onto $\mathbb R$}
 & \GridNo{$\pi_i$ onto $\mathbb R$} \\
\hline
\shortstack[l]{Locally compact\\Hausdorff}
 & \GridYes{Example 3, \S29} & \GridNo{\GridHint{See \cref{prop:sorgenfrey_line_locally_compact}}}
 & \GridNo{\GridHint{See \cref{cor:sorgenfrey_plane_locally_compact}}} & \GridNo{Example 2, \S29}
 & \GridNo{Exercise 4, \S29} & \GridNo{\GridHint{Example 2 \S29}}
 & \GridNo{Example 2, \S29 (adapt)} \\
\hline
Hausdorff
 & \GridYes{Order topology} & \GridYes{\GridHint{$x\neq y \to 
 [x,t), [t, y+1)$}}
 & \GridYes{Thm 31.2(a)} & \GridYes{Thm 31.2(a)}
 & \GridYes{Metrizable} & \GridYes{Refinement of product}
 & \GridYes{Thm 31.2(a)} \\
\hline
Regular
 & \GridYes{$\downarrow$} & \GridYes{$\downarrow$}
 & \GridYes{$\downarrow$} & \GridYes{$\downarrow$}
 & \GridYes{$\downarrow$} & & \GridYes{Thm 31.2(b)} \\
\hline
Completely regular
 & \GridYes{$\downarrow$} & \GridYes{$\downarrow$}
 & \GridYes{Example 1, \S33} & \GridYes{$\downarrow$}
 & \GridYes{$\downarrow$} & & \GridYes{Thm 33.2} \\
\hline
Normal
 & \GridYes{Thm 32.3} & \GridYes{Example 2, \S31}
 & \GridNo{Example 3, \S31} & \GridYes{Thm 32.2}
 & \GridYes{Thm 32.2} & \GridNA
 & \GridNo{Example 1, \S32} \\
\hline
Metrizable
 & \GridWork{\GridNo{See \cref{prop:ordered-square-metrizability}}} & \GridNo{\GridHint{see Exercise 5(a), \S30}}
 & \GridNo{$\leftarrow$} & \GridYes{Thm 20.5}
 & \GridYes{Uniform metric} & \GridNo{\GridHint{}}
 & \GridWork{\GridHint{}} \\
\hline
Second-countable
 & \GridNo{See \cref{prop:ordered_square_second_countability}} & \GridNo{\GridHint{See \cref{prop:sorgenfrey_line_second_countable}}}
 & \GridNo{$\leftarrow$} & \GridYes{Thm 30.2}
 & \GridNo{Example 2, \S30} & \GridNo{Example 2, \S30 (adapt)}
 & \GridWork{\GridHint{$\{0,1\}^{[0,1]}$?}} \\
\hline
\end{tabular}%
}
\end{center}
\endgroup

\bigskip

\begingroup
\setlength{\tabcolsep}{5pt}
\renewcommand{\arraystretch}{1.5}
\begin{center}
\textbf{Exercises 2--6: General preservation properties}

\medskip
\resizebox{\linewidth}{!}{%
\begin{tabular}{|l|*{9}{c|}}
\hline
 & \GridHead{Metric space}
 & \GridHead{Compact Hausdorff}
 & \GridHead{Arbitrary subspace}
 & \GridHead{Open subspace}
 & \GridHead{Closed subspace}
 & \GridHead{Finite products}
 & \GridHead{Countable products}
 & \GridHead{Arbitrary products}
 & \GridHead{Continuous image} \\
\hline
Connected
 & \GridNo{$\{0,1\}_{\mathrm{disc}}$} & \GridNo{$\{0,1\}_{\mathrm{disc}}$}
 & \GridNo{$X\cup Y$} & \GridNo{$X\cup Y$}
 & \GridNo{$\{0,1\}\subset[0,1]$}
 & \GridYes{Thm 23.6} & \GridYes{Exercise 10, \S23}
 & \GridYes{Exercise 10, \S23} & \GridYes{Thm 23.5} \\
\hline
Path connected
 & \GridNo{$\uparrow$} & \GridNo{$\uparrow$}
 & \GridNo{$\mathbb Q \subseteq \mathbb R$} & 
\GridNo{$U \cap V = \emptyset 
\Rightarrow U \cup V$ not connected }
 & \GridNo{Topologists sine curve}  & \GridYes{$\rightarrow$}
 & \GridYes{$\rightarrow$} & \GridYes{See \cref{prop:product_space_path_connected}}& \GridYes{$f\circ p$} \\
\hline
Locally connected
 & \GridNo{$\mathbb Q$} & 
 \GridNo{Topologists sine curve}
 & \GridNo{$\mathbb Q \subseteq \mathbb R$} & 
    \GridYes{$\downarrow$}& \GridNo{Topologists sine curve} & 
 \GridYes{$\downarrow$} & 
  \GridNo{See \cref{prop:product_local_connected}}
 & \GridNo{$\leftarrow$} & 
 \GridNo{$\mathbb R$ discrete $\to$ 
 $\mathbb R$ canonical} \\
\hline
Locally path connected
 & \GridNo{$\uparrow$} & 
 \GridNo{$\uparrow$} & \GridNo{$\mathbb Q \subseteq \mathbb R$} & 
  \GridYes{See \cref{prop:open_subspace_locally_path}}& 
  \GridNo{$\uparrow$}&
 \GridWork{\GridYes{See \cref{prop:product_local_connected}}}
 & \GridNo{$\uparrow$}
 & \GridNo{$\uparrow$} & \GridNo{$\uparrow$} \\
\hline
Compact
 & \GridNo{$\mathbb R$} & \GridYes{By hypothesis}
 & \GridNo{$(0,1)\subset[0,1]$} & \GridNo{$(0,1)\subset[0,1]$}
 & \GridYes{Thm 26.2} & \GridYes{Tychonoff}
 & \GridYes{Tychonoff} & \GridYes{Tychonoff}
 & \GridYes{Thm 26.5} \\
\hline
Locally compact Hausdorff
 & \GridNo{$\uparrow$} & 
 \GridYes{Compact $\rightarrow$ locally compact}
 & & & &
 \GridYes{See \cref{prop:product_local_compact}} & 
 \GridNo{$\mathbb R^\omega$ not 
 locally compact} & \GridNo{$\leftarrow$} & \GridNo{$\downarrow$} \\
\hline
Hausdorff
 & \GridYes{Metric} & \GridYes{By hypothesis}
 & \GridYes{Thm 31.2(a)} & \GridYes{Thm 31.2(a)}
 & \GridYes{Thm 31.2(a)} & \GridYes{Thm 31.2(a)}
 & \GridYes{Thm 31.2(a)} & \GridYes{Thm 31.2(a)}
 & \GridNo{$\{0,1\}$ discrete $\to$ 
 $\{0,1\}$ indiscrete} \\
\hline
Regular
 & \GridYes{$\downarrow$} & \GridYes{$\downarrow$}
 & \GridYes{Thm 31.2(b)} & \GridYes{Thm 31.2(b)}
 & \GridYes{Thm 31.2(b)} & \GridYes{Thm 31.2(b)}
 & \GridYes{Thm 31.2(b)} & \GridYes{Thm 31.2(b)} & 
 \GridNo{See \cref{ex:image_regular}}\\
\hline
Completely regular
 & \GridYes{$\downarrow$} & \GridYes{$\downarrow$}
 & \GridYes{Thm 33.2} & \GridYes{Thm 33.2}
 & \GridYes{Thm 33.2} & \GridYes{Thm 33.2}
 & \GridYes{Thm 33.2} & \GridYes{Thm 33.2} & 
 \GridNo{$\uparrow$} \\
\hline
Normal
 & \GridYes{Thm 32.2} & \GridYes{Thm 32.3}
 & \GridNo{$\rightarrow$} & \GridNo{See 
\cref{ex:open_subspace_normal}}  & 
 \GridYes{Closed in subspace $\Rightarrow$ 
closed in ambient}& \GridNo{$\mathbb R_\ell^2$} & 
 \GridNo{$\leftarrow$}& \GridNo{Example 1, \S32} & 
 \GridNo{Discrete space $\to$ 
 non-normal space} \\
\hline
Second-countable
 & \GridNo{$\mathbb R^\omega$ uniform} & 
 \GridNo{Any uncountable discrete space}
 & \GridYes{Thm 30.2} & \GridYes{Thm 30.2}
 & \GridYes{Thm 30.2} & &
 & \GridNo{$\mathbb R^{[0,1]}$} &  
 \GridNo{$\mathbb R \to \mathbb R/\mathbb Z$ quotient} \\
\hline
Metrizable
 & \GridYes{By definition} & 
 \GridNo{Urysohn Metrization Thm.}
 & \GridYes{Restriction of metric} & \GridYes{Restriction of metric}
 & \GridYes{Restriction of metric} & & & 
 \GridNo{$\mathbb R^{\mathbb R}$ not first countable} & 
 \GridNo{See \cref{ex:image_regular}} \\
\hline
\end{tabular}%
}
\end{center}
\endgroup

\newpage 

\section{Proofs}

\begin{prop}\label{prop:omega_uniform_connectedness}
    $\mathbb R^\omega$ with the uniform topology 
    is \keyterm{not connected}. 
\end{prop}

\begin{proof}
    The countably infinite product 
    of the real line $\mathbb R^\omega$ with 
    the uniform topology is the 
    space induced by the uniform metric 
    \[
        \overline d(x,y) = \sup_{n\in \mathbb N} 
        \min(|x_n-y_n|, 1)
    \]
    where $x = (x_n)$ and $y=(y_n)$ are 
    sequences of real numbers. 

    We claim this space is \keyterm{not connected}. 

    Let $B$ be the set of all bounded sequences, 
    and let $U$ be the set of all unbounded sequences, 
    \begin{enumerate}
        \item $B = \{x\in\mathbb R^\omega \bigm| \exists 
        M > 0 \ \text{s.t.} \ |x_n| \leq M \ \forall n \in 
        \mathbb N\}$
        \item $U = \mathbb R^\omega \setminus B$. 
    \end{enumerate}
    Clearly, these are nonempty since 
    \[
        (0,0,0,\ldots) \in B 
    \]
    while 
    \[
        (1,2,3,\ldots) \in U. 
    \]

    \textbf{$B$ is open.}

    Let $x = (x_n) \in B$. Then there exists 
    some $M_x > 0$ such that $|x_n| \leq M_x$ for all 
    $ n \in \mathbb N$. 

    Consider the ball of radius $1$ around $x$ 
    \[
        B_{\overline d}(x,1) = \{y\in \mathbb R^\omega 
        \bigm| \overline d(x,y) < 1\}.  
    \]
    If $y\in B_{\overline d}(x,1)$, then 
    \[
        \sup_{n\in\mathbb N}\min(|x_n-y_n|,1) < 1
    \]
    implying that for every $n$, $|x_n - y_n| < 1$ for 
    all $n$. 
    Using the triangle inequality we can 
    attain the following bound: 
    \[
        |y_n| \leq |y_n - x_n| + |x_n| < 1 + M_x.
    \]
    Thus the sequence $y$ is bounded, meaning 
    \[
        B_{\overline d}(x,1) \subseteq B. 
    \]
    Since every point of $B$ has an open neighborhood 
    completely contained in $B$, $B$ is open. 
    
    \textbf{$U$ is open.}

    Let $(x_n)$ be undounded. 
    Consider the same open ball $B_{\overline d}(x,1)$. 
    If $y\in B_{\overline d}(x,1)$ then 
    \[
        |y_n| \geq |x_n| - |x_n - y_n| > |x_n| - 1
    \]
    so $y$ is unbounded since $|x_n|$ can be arbitrarily 
    large. 

    Therefore, $B_{\overline d}(x,1) \subseteq U$ is open. 

    Then 
    \[
        \mathbb R^\omega = B \cup U. 
    \]
\end{proof}

\begin{prop}\label{prop:omega_uniform_local_connectedness}
    $\mathbb R^\omega$ with the uniform topology 
    is \keyterm{locally connected}. 
    In fact, it is \keyterm{locally path connected}. 
\end{prop}
\begin{proof}
    We claim the space is 
    \keyterm{locally path connected}. 

    For any point $x = (x_n) \in \mathbb R^\omega$ 
    and any radius $\varepsilon > 0$ the open $\varepsilon$-ball 
    centered at $x$ is 
    \[
        B_{\overline d}(x,\varepsilon)=\{y \in \mathbb R^\omega 
        \bigm| \overline d(x,y) < \varepsilon\}. 
    \]
    If we choose $\varepsilon < 1$, the uniform metric 
    forces every coordinate to be within 
    $\varepsilon$ of $x_n$. 
    Therefore, the open $\varepsilon$-ball 
    is 
    \[
        B_{\overline d}(x,\varepsilon) = 
        \prod_{n\in\mathbb N} (x_n - \varepsilon, 
        x_n + \varepsilon). 
    \]

    Take any point $y \in B_{\overline d}(x,\varepsilon)$. 
    Construct a linear path 
    $f : [0,1] \to B_{\overline d}(x,\varepsilon)$ 
    from the center point $x$ to $y$ by setting 
    \[
        f(t) = (1-t)x+ty. 
    \]
    Since each interval $(x_n - \varepsilon, 
    x_n + \varepsilon)$ is convex, 
    every point remains inside the ball 
    for all $t \in [0,1]$. 

    Since every point has a neighborhood basis 
    consisting of these path-connected open balls, 
    the space is locally path-connected which makes 
    it locally connected. 

\end{proof}

\begin{prop}\label{prop:ordered-square-metrizability}
    The ordered square $I^2_o$ is 
    \keyterm{not metrizable}. 
\end{prop}
\begin{proof}[Proof sketch]
    We invoke the following theorem: 
    \begin{thm*}
        If a topological space is \keyterm{compact} 
        and \keyterm{metrizable}, 
        then it must be separable. 
    \end{thm*}

    So it will suffice to elucidate that the space 
    is compact and not separable. 

    The Heine-Borel theorem for ordered spaces says that 
    a linearly ordered topological space is compact if 
    and only if it has a least element, a greatest 
    element, and satisfies the least upper 
    bound properrty. 

    The ordered square has minimum $(0,0)$, 
    maximum $(1,1)$. 
    Because the standard unit interval has 
    the least upper bound property, the lexicographic 
    order preserves this completeness. 
    Thus $I_o^2$ is compact. 

    The space is not separable. 
    We will find an uncountable 
    collection of open sets that are completely 
    isolated from each other. 
    For every real $x \in (0,1)$ consider the 
    vertical segment at $x$, excluding endpoints, 
    \[
        V_x = \{x\} \times (0,1). 
    \]
    In the dictionary order, $V_x$ is the 
    open interval $((x,0), (x,1))$. Because it is an 
    open interval, $V_x$ is an open set in the dictionary 
    order topology. 
    Notice that if $x \neq y$ the vertical segments 
    $V_x$ and $V_y$ are disjoint, 
    \[
        V_x \cap V_y = \emptyset. 
    \]
    Since there are uncountably many choices 
    of $x \in (0,1)$ we have an uncountable 
    collection of pairwise disjoint non-empty 
    open sets. 

    Suppose then that some 
    dense $D \subseteq I_o^2$ exists. 
    By definition the set must intersect 
    all non-empty open sets, but there are uncountably 
    many so this cannot be the case. 

\end{proof}

\begin{prop}\label{prop:ordered_square_second_countability}
    The ordered square $I_o^2$ is 
    \keyterm{not second-countable}.
\end{prop}
\begin{proof}[Proof sketch]
    As we saw in the proof of 
    \cref{prop:ordered-square-metrizability}, 
    we can construct an uncountable 
    collection of pairwise disjoint, non-empty 
    open sets. 
    By definition, each of these 
    sets must contain at least 
    one basic open set from the basis. 

    Since the vertical intervals are disjoint 
    from one another, a basic open set 
    in one cannot be shared by any other, 
    meaning we must have uncountably many basis elements. 
\end{proof}

\begin{prop}\label{prop:sorgenfrey_line_locally_compact}
    The Sorgenfrey line $\mathbb R_\ell$ is 
    \keyterm{not locally compact Hausdorff}. 
\end{prop}
\begin{proof}[Proof by some guy on mathstackexchange said]
    In this 
    \href{https://math.stackexchange.com/questions/169896/how-to-prove-that-compact-subspaces-of-the-sorgenfrey-line-are-countable/170709#170709}{Math stackexchange post} there is given a proof that 
    the Sorgenfrey line is not locally compact, and thus 
    not locally compact Hausdorff. 
\end{proof}

\begin{cor}\label{cor:sorgenfrey_plane_locally_compact}
    The Sorgenfrey plane $\mathbb R^2_\ell$ is 
    \keyterm{not locally compact Hausdorff}. 
\end{cor}
\begin{proof}[Proof sketch]
    Suppose it was. 
    Considert the standard projection map 
    onto coordinate one, 
    \[
        \pi_1 : \mathbb R_\ell \times \mathbb R_\ell 
        \to \mathbb R_\ell
    \]
    where 
    \[
        \pi_1(x,y) = x.  
    \]
    This map is continuous and open. 
    The image of a locally compact space 
    under a continuous open map is also 
    locally compact, meaning $\mathbb R_\ell$ is locally compact, 
    contradiction \cref{prop:sorgenfrey_line_locally_compact}. 

    Thus $\mathbb R^2_\ell$ is not locally compact and thus 
    not locally compact Hausdorff. 
\end{proof}

\begin{prop}\label{prop:sorgenfrey_line_second_countable}
    The Sorgenfrey line $\mathbb R_\ell$ is \keyterm{not 
    second countable}.
\end{prop}
\begin{proof}
    Suppose it is second-countable with 
    basis 
    \[
        \mathcal B = \{B_1,B_2,\ldots\}.
    \]

    For every $x \in \mathbb R$, 
    $[x,x+1)$ is open. 
    Since $\mathcal B$ is a basis there must exist, 
    for every real $x \in \mathbb R$, an element $B_{n_x} \in 
    \mathcal B$ such that 
    \[
        x \in B_{n_x} \subseteq [x,x+1]. 
    \]
    Notice that $x$ must be the absolute least element 
    in $[x,x+1)$. Since $B_{n_x} \subseteq [x,x+1]$ and 
    $x \in B_{n_x}$, $x$ must be the absolute least element 
    of $B_{n_x}$. 
    
    If we have two distinct $x,y \in \mathbb R$, then 
    they cannot share the same basis element. 
    If $B_{n_x} = B_{n_y}$ then these 
    basis elements have two distinct minimum values, 
    which is impossible. 

    Thus we can construct an injection 
    $x \mapsto B_{n_x}$, meaning 
    \[
        \aleph_1 = |\mathbb R_\ell| \leq |\mathcal B| 
    \]
    so $\mathcal B$ is uncountable. 
\end{proof}

\begin{prop}\label{prop:box_locally_path_connected}
    $\mathbb R^\omega$ with the box topology 
    is \keyterm{not locally path connected}. 
\end{prop}
\begin{proof}
    Let $U$ be a basic open neighborhood of $0$ 
    by 
    \[
        U = \prod_{n\in\mathbb N}(-1/n, 1/n). 
    \]
    Pick a point $x = (x_1,x_2,\ldots) \in \mathbb R^\omega$ 
    by setting 
    \[
        x_n = \frac{1}{2n}. 
    \]
    Since $x_n \in (-1/n, 1/n)$ for all $n \in \mathbb N$ 
    we have $x \in U$. 

    Suppose a path $f : [0,1] \to U$ with $f(0) = 0, f(1) = x$ 
    exists. 

    Then 
    \[
        f_n = \pi_n \circ f : [0,1] \to (-1/n, 1/n)
    \]
    is a continuous real valued function with 
    $f_n(0) = 0$ and $f_n(1) = 1/2n$. 

    By the intermediate value theorem, each 
    $f_n$ attains all values between $0$ and $1/2n$. 

    In particular, there exists $t_n \in [0,1]$ so 
    \[
        f_n(t_n) = \frac{1}{4n}. 
    \]

    We get a sequence $\{t_n\}_n \subseteq [0,1]$, 
    a compact metric space. Let $t^* \in [0,1]$ be the limit of 
    some convergent subsequence. Assume without loss 
    of generality that 
    \[
        \lim t_n = t^*.
    \]

    Let $f(t^*) = (y_1,y_2,\ldots) \in U$. 
    Take 
    \[
        V = \prod_{n\in\mathbb N} (y_n - \varepsilon_n, 
        y_n + \varepsilon_n)
    \]
    where 
    \[
        \varepsilon_n = \frac{1}{2}\left|
            \frac{1}{4n} - y_n
        \right| < |f_n(t_n) - y_n|. 
    \]
    If $y_n = 1/4n$ set $\varepsilon_n = 1/8n$.

    Then 
    \[
        \left|\frac{1}{4n} - y_n\right| \geq \varepsilon_n > 0.
    \]
    Since $V$ is open, there exists $\delta > 0$ so 
    \[
        |t-t^*| < \delta \Rightarrow f(t) \in V. 
    \]
    Since $\lim t_n = t^*$ choose $N$ large enough so 
    \[
        |t_n - t^*| < \delta, \ n \geq N. 
    \]
    Then $f(t_n) \in V$.

    Now, for all $n \geq N$ we have 
    \[
        \pi_n(f(t_n)) \in (y_n - \varepsilon_n, y_n 
        + \varepsilon_n)
    \]
    meaning 
    \[
        |f_n(t_n) - y_n| < \varepsilon_n. 
    \]
    Substituting we get 
    \[
        \left|\frac{1}{4n} - y_n\right| < \varepsilon_n, 
    \]
    but this contradicts 
    \[
        \left|\frac{1}{4n} - y_n\right| \geq \varepsilon_n.
    \]
\end{proof}

\begin{prop}\label{prop:product_space_path_connected}
    If $X_i$, ranging 
    over indices $i \in I$ are all path connected. 
    Then 
    \[
        X := \prod_{i\in I} X_i 
    \]
    \keyterm{is path connected} (with the product topology). 

    Note: This fails quite spectacularly if we 
    assume the box topology. 
\end{prop}
\begin{proof}[Proof sketch]
    Recall that the basis is 
    \[
        \mathcal B = \left\{
            \prod_{i \in I} U_i \Biggm| 
            U_i \ \text{open} \ \text{in} \ X_i \ \text{and} \ 
            U_i \neq X_i \ \text{for finitely many} \ i \in I
        \right\}.
    \]

    Take some $x = (x_i)_{i\in I}, \ y = 
    (y_i)_{i \in I} \in X$. 

    Then, since $X_i$ is path connected 
    for each $i \in I$, we have 
    some continuous 
    $f_{i} : [0,1] \to X_i$ with 
    \[
        f(0) = x_i, \ \text{and} \ f(1) = y_i. 
    \]

    Construct the path $f : [0,1] \to X$ by 
    \[
        f(t) = (f_i(t))_{i \in I}.   
    \]

    This is nice and well-defined because I say so and we have 
    \[
        f(0) = (f_i(0))_{i \in I} 
        = (x_i)_{i\in I}
    \]
    and 
    \[
        f(1) = (f_i(1))_{i \in I} 
        = (y_i)_{i \in I}. 
    \]

    By the universal property of the 
    product topology, $f$ is continuous if and only if 
    $\pi_i \circ f = f_i$ is continuous for all 
    $i \in I$. 
    
    Therefore $f$ is continuous.
\end{proof}

\begin{prop}\label{prop:product_local_connected}
    Finite products of locally connected spaces 
    \keyterm{are locally path connected}, 
    whereas countably infinite products, and thus 
    arbitrary products, of locally connected spaces 
    fail to be even locally connected.  
\end{prop}
\begin{proof}
    Suppose $X_i$, with $i \in [n] := \{1,2,\ldots,n\}$, are 
    locally path connected. 

    Set 
    \[
        X = \prod_{i\in[n]} X_i.
    \]

    A sub-basis can be written as 
    $\pi_i^{-1}(U_i)$ where $u_i \subseteq X_i$ is open. 
    
    A basic open set $W \subset X$ containing 
    some point $x = (x_1,x_2,\ldots)$ can be written 
    as 
    \[
        W = \prod_{i\in[n]} U_i
    \]
    where $U_i$ is an open neighborhood of $x_i$. 

    Since $X_i$ is locally path connected, 
    for each $x_i \in U_i$ we can find 
    a path connected open neighborhood $V_i \subset X_i$ 
    such that 
    \[
        x_i \in V_i \subseteq U_i. 
    \]
    Then 
    \[
        V = V_1\times V_2 \times \cdots \times V_n 
    \]
    is an open subset of $X$ in the product topology. 

    It is path connected since, as has been established, 
    the product of path connected spaces is path connected. 

    Thus $X$ is locally path connected.

    \keyterm{Counterexample for countably infinite product.}

    Take $X = \{0,1\}^{\mathbb N}$, i.e. the Cantor space 
    with the product topology. 

    Each $X_n = \{0,1\}$ are trivially locally path connected 
    and thus locally connected. 

    Let $C \subseteq X$ be connected. 

    For each $n \in \mathbb N$ consider the 
    restriction 
    \[
        \pi_n\bigm|_C : C \to \{0,1\}. 
    \]
    Since $\pi_n$ is continuous and 
    $C$ is connected, the image must be, but 
    $\{0,1\}$ is discrete so its onoly nonempty 
    connected subsets are singletons.- 

    Consequently, for each $n$, there exists 
    $c_n = 0,1$ such that 
    \[
        \pi_n(C) = \{c_n\}. 
    \]

    Thus every element of $C$ has the 
    same coordinates, so 
    \[
        C = \{(c_n)_{n\in \mathbb N}\}. 
    \]
    Thus $X$ is totally disconnected. 

\end{proof}

\begin{ex}\label{ex:image_regular}
    \medskip 
    \ 
    
    \keyterm{regular $\not\rightarrow$ regular \& 
    metrizable $\not\rightarrow$ metrizable}

    Let $X = \mathbb R$ with the 
    standard topology, a metric space 
    and thus regular. Furtheremore the space is metrizable.

    Let $Y = \{0,1\}$ with the trivial topology. 

    Let $f : \mathbb R \to \{0,1\}$ with 
    \[
        f(x) = 
        \begin{cases}
            0 & \text{if} \ x \in \mathbb Q \\ 
            1 & \text{if} \ x \in \mathbb R \setminus 
            \mathbb Q. 
        \end{cases}
    \]

    We have $f^{-1}(Y) = \mathbb R$, open, and 
    $f^{-1}(\emptyset) = \emptyset$, open. 

    $Y$ is not regular since it is not $T-1$, 
    $\{0\} = Y \setminus \{1\}$ and $\{1\}$ is not 
    open. 

    Furthermore, $Y$ is not metrizable since it is not 
    Hausdorff.

\end{ex}

\begin{prop}\label{prop:product_local_compact}
    The finite product of \keyterm{locally compact} spaces
    is \keyterm{locally compact}.
\end{prop}
\begin{proof}[Proof sketch]
    Let $X_i$ for $i = 1,2,3,\ldots,n$ be 
    locally compact. 

    Then every point in each $X_i$ has a compact 
    neighborhood. 

    Take any point $x = (x_1,x_2,\ldots)$. 
    For each $x_i \in X_i$ choose compact neighborhood 
    $K_i \subseteq X_i$. 

    The cartesian product 
    \[
        K = \prod_{i\in[n]} K_i 
    \]
    is compact (apply Tychonoff's theorem or 
    go directly from the fact that the product is 
    finite). 

    Moreover, $K$ is a neighborhood of $x$ in the 
    product topology. 

    Thus, every point has a compact neighborhood, making 
    the finite product locally compact. 
\end{proof}

\begin{prop}\label{prop:open_subspace_locally_path}
    An open subspace of a locally path connected space is locally path 
    connected. 
\end{prop}
\begin{proof}[Proof sketch]
    Let $X$ be locally path connected and 
    $Y \subseteq Y$ open. 

    Let $x \in Y$ and $U \subseteq Y$ and open neighborhood of 
    $x$. Then, since $Y$ is open, $U$ is open in $X$. 
    Thus there exists $W_U \subseteq U$ path connected such that 
    \[
        x \in W_U \subseteq U, 
    \]
    but 
    \[
        U \subseteq Y 
    \]
    so 
    \[
        W_U \subseteq Y 
    \]
    is open path connected in $Y$. 
\end{proof}

\begin{ex}\label{ex:open_subspace_normal}
    Let $X = \{a,b,c,d\}$ with 
    $\mathcal T = \{\emptyset, X, \{d\}, \{b,d\}, 
    \{c,d\}, \{b,c,d\}\}$. 
    $X$ is normal as there 
    exist no disjoint nonempty closed sets (they all 
    have $a$ in common). 

    The space is normal, but 
    $Y = \{b,c,d\}$ (open) is not 
    since $\{b\}, \{d\}$ cannot be separated 
    into disjoint open sets, 
    \[
        \mathcal T_Y \setminus \{\emptyset, Y\} = \{
        \{d\}, \{b,d\}, \{c,d\}\}. 
    \]
\end{ex}


\end{document}
